\documentclass[letterpaper,12pt]{article}
\usepackage[margin=.5in]{geometry}
\usepackage{tikz}
\usepackage{helvet}
\renewcommand{\familydefault}{\sfdefault}
\pagestyle{empty}
\pdfmapfile{+cm.map}
\pdfmapfile{+ps2pk35.map}
\setlength{\parindent}{0pt}
\hyphenpenalty=10000
\exhyphenpenalty=10000
\begin{document}
\null\begin{tikzpicture}[remember picture,overlay,x=1cm,y=1cm]
\begin{scope}[shift={(current page.north west)}]
\node[anchor=north west,inner sep=0pt,font=\fontsize{11.5}{13}\selectfont] at (1.4,-.72) {Week 31 / Hidden orchard / Grades 4--5};
\node[anchor=south west,inner sep=0pt,font=\fontsize{9}{11}\selectfont] at (1.4,-27.2) {Bellingham Math Circle / Week 31 / F31-45-v2};
\node[anchor=south east,inner sep=0pt,font=\fontsize{9}{11}\selectfont] at (20.15,-27.2) {1};
\begin{scope}[shift={(1.4,-1.65)}]
\node[anchor=north west,inner sep=0pt,align=left,text width=18.5cm,font=\fontsize{13}{16.38}\selectfont] at (0,0) {Stretch a thread from O to another dot T. A dot exactly on the thread between O and T blocks T. B marks a blocker. Keep O fixed while you test.};
\draw[black,line width=.55pt] (1,-6.2) -- (4.6,-4.4);
\fill (1.0,-6.2) circle (1.3pt);
\fill (1.0,-5.300000000000001) circle (1.3pt);
\fill (1.0,-4.4) circle (1.3pt);
\fill (1.0,-3.5) circle (1.3pt);
\fill (1.0,-2.6) circle (1.3pt);
\fill (1.9,-6.2) circle (1.3pt);
\fill (1.9,-5.300000000000001) circle (1.3pt);
\fill (1.9,-4.4) circle (1.3pt);
\fill (1.9,-3.5) circle (1.3pt);
\fill (1.9,-2.6) circle (1.3pt);
\fill (2.8,-6.2) circle (1.3pt);
\fill (2.8,-5.300000000000001) circle (1.3pt);
\fill (2.8,-4.4) circle (1.3pt);
\fill (2.8,-3.5) circle (1.3pt);
\fill (2.8,-2.6) circle (1.3pt);
\fill (3.7,-6.2) circle (1.3pt);
\fill (3.7,-5.300000000000001) circle (1.3pt);
\fill (3.7,-4.4) circle (1.3pt);
\fill (3.7,-3.5) circle (1.3pt);
\fill (3.7,-2.6) circle (1.3pt);
\fill (4.6,-6.2) circle (1.3pt);
\fill (4.6,-5.300000000000001) circle (1.3pt);
\fill (4.6,-4.4) circle (1.3pt);
\fill (4.6,-3.5) circle (1.3pt);
\fill (4.6,-2.6) circle (1.3pt);
\draw[line width=.8pt] (4.6,-4.4) circle (3.7pt);
\draw[line width=.8pt] (2.8,-5.300000000000001) circle (3.7pt);
\node[font=\fontsize{10}{12}\selectfont,inner sep=1pt] at (1.0,-6.68) {0};
\node[font=\fontsize{10}{12}\selectfont,inner sep=1pt] at (1.9,-6.68) {1};
\node[font=\fontsize{10}{12}\selectfont,inner sep=1pt] at (2.8,-6.68) {2};
\node[font=\fontsize{10}{12}\selectfont,inner sep=1pt] at (3.7,-6.68) {3};
\node[font=\fontsize{10}{12}\selectfont,inner sep=1pt] at (4.6,-6.68) {4};
\node[font=\fontsize{10}{12}\selectfont,inner sep=1pt] at (0.5,-5.300000000000001) {1};
\node[font=\fontsize{10}{12}\selectfont,inner sep=1pt] at (0.5,-4.4) {2};
\node[font=\fontsize{10}{12}\selectfont,inner sep=1pt] at (0.5,-3.5) {3};
\node[font=\fontsize{10}{12}\selectfont,inner sep=1pt] at (0.5,-2.6) {4};
\node[font=\fontsize{12}{14}\selectfont,inner sep=1pt] at (0.55,-6.25) {O};
\node[font=\fontsize{12}{14}\selectfont,inner sep=1pt] at (4.9,-4.15) {T};
\node[font=\fontsize{12}{14}\selectfont,inner sep=1pt] at (2.4,-5.3) {B};
\node[font=\fontsize{12}{14}\selectfont,inner sep=1pt] at (2.8,-7) {hidden};
\draw[black,line width=.55pt] (10,-6.2) -- (12.7,-4.4);
\fill (10.0,-6.2) circle (1.3pt);
\fill (10.0,-5.300000000000001) circle (1.3pt);
\fill (10.0,-4.4) circle (1.3pt);
\fill (10.0,-3.5) circle (1.3pt);
\fill (10.0,-2.6) circle (1.3pt);
\fill (10.9,-6.2) circle (1.3pt);
\fill (10.9,-5.300000000000001) circle (1.3pt);
\fill (10.9,-4.4) circle (1.3pt);
\fill (10.9,-3.5) circle (1.3pt);
\fill (10.9,-2.6) circle (1.3pt);
\fill (11.8,-6.2) circle (1.3pt);
\fill (11.8,-5.300000000000001) circle (1.3pt);
\fill (11.8,-4.4) circle (1.3pt);
\fill (11.8,-3.5) circle (1.3pt);
\fill (11.8,-2.6) circle (1.3pt);
\fill (12.7,-6.2) circle (1.3pt);
\fill (12.7,-5.300000000000001) circle (1.3pt);
\fill (12.7,-4.4) circle (1.3pt);
\fill (12.7,-3.5) circle (1.3pt);
\fill (12.7,-2.6) circle (1.3pt);
\fill (13.6,-6.2) circle (1.3pt);
\fill (13.6,-5.300000000000001) circle (1.3pt);
\fill (13.6,-4.4) circle (1.3pt);
\fill (13.6,-3.5) circle (1.3pt);
\fill (13.6,-2.6) circle (1.3pt);
\draw[line width=.8pt] (12.7,-4.4) circle (3.7pt);
\node[font=\fontsize{10}{12}\selectfont,inner sep=1pt] at (10.0,-6.68) {0};
\node[font=\fontsize{10}{12}\selectfont,inner sep=1pt] at (10.9,-6.68) {1};
\node[font=\fontsize{10}{12}\selectfont,inner sep=1pt] at (11.8,-6.68) {2};
\node[font=\fontsize{10}{12}\selectfont,inner sep=1pt] at (12.7,-6.68) {3};
\node[font=\fontsize{10}{12}\selectfont,inner sep=1pt] at (13.6,-6.68) {4};
\node[font=\fontsize{10}{12}\selectfont,inner sep=1pt] at (9.5,-5.300000000000001) {1};
\node[font=\fontsize{10}{12}\selectfont,inner sep=1pt] at (9.5,-4.4) {2};
\node[font=\fontsize{10}{12}\selectfont,inner sep=1pt] at (9.5,-3.5) {3};
\node[font=\fontsize{10}{12}\selectfont,inner sep=1pt] at (9.5,-2.6) {4};
\node[font=\fontsize{12}{14}\selectfont,inner sep=1pt] at (9.55,-6.25) {O};
\node[font=\fontsize{12}{14}\selectfont,inner sep=1pt] at (13.05,-4.15) {T};
\node[font=\fontsize{12}{14}\selectfont,inner sep=1pt] at (11.8,-7) {visible};
\node[anchor=north west,inner sep=0pt,align=left,text width=18.6cm,font=\fontsize{14}{17.78}\selectfont] at (0,-8.4) {\textbf{Problem 1:} Find three visible targets and three hidden targets with different numbers of blockers. Mark each blocker.};
\fill (3,-23) circle (1.3pt);
\fill (3,-21) circle (1.3pt);
\fill (3,-19) circle (1.3pt);
\fill (3,-17) circle (1.3pt);
\fill (3,-15) circle (1.3pt);
\fill (3,-13) circle (1.3pt);
\fill (3,-11) circle (1.3pt);
\fill (5,-23) circle (1.3pt);
\fill (5,-21) circle (1.3pt);
\fill (5,-19) circle (1.3pt);
\fill (5,-17) circle (1.3pt);
\fill (5,-15) circle (1.3pt);
\fill (5,-13) circle (1.3pt);
\fill (5,-11) circle (1.3pt);
\fill (7,-23) circle (1.3pt);
\fill (7,-21) circle (1.3pt);
\fill (7,-19) circle (1.3pt);
\fill (7,-17) circle (1.3pt);
\fill (7,-15) circle (1.3pt);
\fill (7,-13) circle (1.3pt);
\fill (7,-11) circle (1.3pt);
\fill (9,-23) circle (1.3pt);
\fill (9,-21) circle (1.3pt);
\fill (9,-19) circle (1.3pt);
\fill (9,-17) circle (1.3pt);
\fill (9,-15) circle (1.3pt);
\fill (9,-13) circle (1.3pt);
\fill (9,-11) circle (1.3pt);
\fill (11,-23) circle (1.3pt);
\fill (11,-21) circle (1.3pt);
\fill (11,-19) circle (1.3pt);
\fill (11,-17) circle (1.3pt);
\fill (11,-15) circle (1.3pt);
\fill (11,-13) circle (1.3pt);
\fill (11,-11) circle (1.3pt);
\fill (13,-23) circle (1.3pt);
\fill (13,-21) circle (1.3pt);
\fill (13,-19) circle (1.3pt);
\fill (13,-17) circle (1.3pt);
\fill (13,-15) circle (1.3pt);
\fill (13,-13) circle (1.3pt);
\fill (13,-11) circle (1.3pt);
\fill (15,-23) circle (1.3pt);
\fill (15,-21) circle (1.3pt);
\fill (15,-19) circle (1.3pt);
\fill (15,-17) circle (1.3pt);
\fill (15,-15) circle (1.3pt);
\fill (15,-13) circle (1.3pt);
\fill (15,-11) circle (1.3pt);
\node[font=\fontsize{10}{12}\selectfont,inner sep=1pt] at (3,-23.48) {0};
\node[font=\fontsize{10}{12}\selectfont,inner sep=1pt] at (5,-23.48) {1};
\node[font=\fontsize{10}{12}\selectfont,inner sep=1pt] at (7,-23.48) {2};
\node[font=\fontsize{10}{12}\selectfont,inner sep=1pt] at (9,-23.48) {3};
\node[font=\fontsize{10}{12}\selectfont,inner sep=1pt] at (11,-23.48) {4};
\node[font=\fontsize{10}{12}\selectfont,inner sep=1pt] at (13,-23.48) {5};
\node[font=\fontsize{10}{12}\selectfont,inner sep=1pt] at (15,-23.48) {6};
\node[font=\fontsize{10}{12}\selectfont,inner sep=1pt] at (2.5,-21) {1};
\node[font=\fontsize{10}{12}\selectfont,inner sep=1pt] at (2.5,-19) {2};
\node[font=\fontsize{10}{12}\selectfont,inner sep=1pt] at (2.5,-17) {3};
\node[font=\fontsize{10}{12}\selectfont,inner sep=1pt] at (2.5,-15) {4};
\node[font=\fontsize{10}{12}\selectfont,inner sep=1pt] at (2.5,-13) {5};
\node[font=\fontsize{10}{12}\selectfont,inner sep=1pt] at (2.5,-11) {6};
\node[font=\fontsize{12}{14}\selectfont,inner sep=1pt] at (2.55,-23.05) {O};\end{scope}\end{scope}\end{tikzpicture}
\newpage
\null\begin{tikzpicture}[remember picture,overlay,x=1cm,y=1cm]
\begin{scope}[shift={(current page.north west)}]
\node[anchor=north west,inner sep=0pt,font=\fontsize{11.5}{13}\selectfont] at (1.4,-.72) {Week 31 / Hidden orchard / Grades 4--5};
\node[anchor=south west,inner sep=0pt,font=\fontsize{9}{11}\selectfont] at (1.4,-27.2) {Bellingham Math Circle / Week 31 / F31-45-v2};
\node[anchor=south east,inner sep=0pt,font=\fontsize{9}{11}\selectfont] at (20.15,-27.2) {2};
\begin{scope}[shift={(1.4,-1.65)}]
\node[anchor=north west,inner sep=0pt,align=left,text width=18.5cm,font=\fontsize{13}{16.38}\selectfont] at (0,0) {(4, 2) means 4 across and 2 up from O.};
\fill (1.0,-4.1) circle (1.3pt);
\fill (1.0,-3.3999999999999995) circle (1.3pt);
\fill (1.0,-2.7) circle (1.3pt);
\fill (1.0,-2.0) circle (1.3pt);
\fill (1.0,-1.3) circle (1.3pt);
\fill (1.7,-4.1) circle (1.3pt);
\fill (1.7,-3.3999999999999995) circle (1.3pt);
\fill (1.7,-2.7) circle (1.3pt);
\fill (1.7,-2.0) circle (1.3pt);
\fill (1.7,-1.3) circle (1.3pt);
\fill (2.4,-4.1) circle (1.3pt);
\fill (2.4,-3.3999999999999995) circle (1.3pt);
\fill (2.4,-2.7) circle (1.3pt);
\fill (2.4,-2.0) circle (1.3pt);
\fill (2.4,-1.3) circle (1.3pt);
\fill (3.0999999999999996,-4.1) circle (1.3pt);
\fill (3.0999999999999996,-3.3999999999999995) circle (1.3pt);
\fill (3.0999999999999996,-2.7) circle (1.3pt);
\fill (3.0999999999999996,-2.0) circle (1.3pt);
\fill (3.0999999999999996,-1.3) circle (1.3pt);
\fill (3.8,-4.1) circle (1.3pt);
\fill (3.8,-3.3999999999999995) circle (1.3pt);
\fill (3.8,-2.7) circle (1.3pt);
\fill (3.8,-2.0) circle (1.3pt);
\fill (3.8,-1.3) circle (1.3pt);
\draw[line width=.8pt] (3.8,-2.7) circle (3.7pt);
\node[font=\fontsize{10}{12}\selectfont,inner sep=1pt] at (1.0,-4.58) {0};
\node[font=\fontsize{10}{12}\selectfont,inner sep=1pt] at (1.7,-4.58) {1};
\node[font=\fontsize{10}{12}\selectfont,inner sep=1pt] at (2.4,-4.58) {2};
\node[font=\fontsize{10}{12}\selectfont,inner sep=1pt] at (3.0999999999999996,-4.58) {3};
\node[font=\fontsize{10}{12}\selectfont,inner sep=1pt] at (3.8,-4.58) {4};
\node[font=\fontsize{10}{12}\selectfont,inner sep=1pt] at (0.5,-3.3999999999999995) {1};
\node[font=\fontsize{10}{12}\selectfont,inner sep=1pt] at (0.5,-2.7) {2};
\node[font=\fontsize{10}{12}\selectfont,inner sep=1pt] at (0.5,-2.0) {3};
\node[font=\fontsize{10}{12}\selectfont,inner sep=1pt] at (0.5,-1.3) {4};
\node[font=\fontsize{12}{14}\selectfont,inner sep=1pt] at (0.55,-4.1499999999999995) {O};
\draw[->,line width=.9pt] (1,-4.1) -- (3.8,-4.1);
\draw[->,line width=.9pt] (3.8,-4.1) -- (3.8,-2.7);
\node[font=\fontsize{11}{13}\selectfont,inner sep=1pt] at (2.4,-4.95) {4 across};
\node[font=\fontsize{11}{13}\selectfont,inner sep=1pt] at (5.0,-3.4) {2 up};
\node[font=\fontsize{12}{14}\selectfont,inner sep=1pt] at (5.9,-2.15) {(4, 2)};
\node[anchor=north west,inner sep=0pt,align=left,text width=18.6cm,font=\fontsize{14}{17.78}\selectfont] at (0,-6.3) {\textbf{Problem 2:} For each target, find the first dot after O on its line and every blocker: (6, 4), (6, 3), (5, 3), (6, 6), (6, 0) and (0, 6).};
\fill (3,-21.5) circle (1.3pt);
\fill (3,-19.5) circle (1.3pt);
\fill (3,-17.5) circle (1.3pt);
\fill (3,-15.5) circle (1.3pt);
\fill (3,-13.5) circle (1.3pt);
\fill (3,-11.5) circle (1.3pt);
\fill (3,-9.5) circle (1.3pt);
\fill (5,-21.5) circle (1.3pt);
\fill (5,-19.5) circle (1.3pt);
\fill (5,-17.5) circle (1.3pt);
\fill (5,-15.5) circle (1.3pt);
\fill (5,-13.5) circle (1.3pt);
\fill (5,-11.5) circle (1.3pt);
\fill (5,-9.5) circle (1.3pt);
\fill (7,-21.5) circle (1.3pt);
\fill (7,-19.5) circle (1.3pt);
\fill (7,-17.5) circle (1.3pt);
\fill (7,-15.5) circle (1.3pt);
\fill (7,-13.5) circle (1.3pt);
\fill (7,-11.5) circle (1.3pt);
\fill (7,-9.5) circle (1.3pt);
\fill (9,-21.5) circle (1.3pt);
\fill (9,-19.5) circle (1.3pt);
\fill (9,-17.5) circle (1.3pt);
\fill (9,-15.5) circle (1.3pt);
\fill (9,-13.5) circle (1.3pt);
\fill (9,-11.5) circle (1.3pt);
\fill (9,-9.5) circle (1.3pt);
\fill (11,-21.5) circle (1.3pt);
\fill (11,-19.5) circle (1.3pt);
\fill (11,-17.5) circle (1.3pt);
\fill (11,-15.5) circle (1.3pt);
\fill (11,-13.5) circle (1.3pt);
\fill (11,-11.5) circle (1.3pt);
\fill (11,-9.5) circle (1.3pt);
\fill (13,-21.5) circle (1.3pt);
\fill (13,-19.5) circle (1.3pt);
\fill (13,-17.5) circle (1.3pt);
\fill (13,-15.5) circle (1.3pt);
\fill (13,-13.5) circle (1.3pt);
\fill (13,-11.5) circle (1.3pt);
\fill (13,-9.5) circle (1.3pt);
\fill (15,-21.5) circle (1.3pt);
\fill (15,-19.5) circle (1.3pt);
\fill (15,-17.5) circle (1.3pt);
\fill (15,-15.5) circle (1.3pt);
\fill (15,-13.5) circle (1.3pt);
\fill (15,-11.5) circle (1.3pt);
\fill (15,-9.5) circle (1.3pt);
\node[font=\fontsize{10}{12}\selectfont,inner sep=1pt] at (3,-21.98) {0};
\node[font=\fontsize{10}{12}\selectfont,inner sep=1pt] at (5,-21.98) {1};
\node[font=\fontsize{10}{12}\selectfont,inner sep=1pt] at (7,-21.98) {2};
\node[font=\fontsize{10}{12}\selectfont,inner sep=1pt] at (9,-21.98) {3};
\node[font=\fontsize{10}{12}\selectfont,inner sep=1pt] at (11,-21.98) {4};
\node[font=\fontsize{10}{12}\selectfont,inner sep=1pt] at (13,-21.98) {5};
\node[font=\fontsize{10}{12}\selectfont,inner sep=1pt] at (15,-21.98) {6};
\node[font=\fontsize{10}{12}\selectfont,inner sep=1pt] at (2.5,-19.5) {1};
\node[font=\fontsize{10}{12}\selectfont,inner sep=1pt] at (2.5,-17.5) {2};
\node[font=\fontsize{10}{12}\selectfont,inner sep=1pt] at (2.5,-15.5) {3};
\node[font=\fontsize{10}{12}\selectfont,inner sep=1pt] at (2.5,-13.5) {4};
\node[font=\fontsize{10}{12}\selectfont,inner sep=1pt] at (2.5,-11.5) {5};
\node[font=\fontsize{10}{12}\selectfont,inner sep=1pt] at (2.5,-9.5) {6};
\node[font=\fontsize{12}{14}\selectfont,inner sep=1pt] at (2.55,-21.55) {O};\end{scope}\end{scope}\end{tikzpicture}
\newpage
\null\begin{tikzpicture}[remember picture,overlay,x=1cm,y=1cm]
\begin{scope}[shift={(current page.north west)}]
\node[anchor=north west,inner sep=0pt,font=\fontsize{11.5}{13}\selectfont] at (1.4,-.72) {Week 31 / Hidden orchard / Grades 4--5};
\node[anchor=south west,inner sep=0pt,font=\fontsize{9}{11}\selectfont] at (1.4,-27.2) {Bellingham Math Circle / Week 31 / F31-45-v2};
\node[anchor=south east,inner sep=0pt,font=\fontsize{9}{11}\selectfont] at (20.15,-27.2) {3};
\begin{scope}[shift={(1.4,-1.65)}]
\node[anchor=north west,inner sep=0pt,align=left,text width=18.6cm,font=\fontsize{14}{17.78}\selectfont] at (0,0) {\textbf{Problem 3:} Predict the first dot and number of blockers for (12, 8), (15, 10), (21, 14), (25, 15) and (13, 8). Explain a rule that works for any target with whole-number coordinates.};
\draw[gray!45,line width=.35pt] (0,-3.0) -- (18.5,-3.0);
\draw[gray!45,line width=.35pt] (0,-4.15) -- (18.5,-4.15);
\draw[gray!45,line width=.35pt] (0,-5.3) -- (18.5,-5.3);
\draw[gray!45,line width=.35pt] (0,-6.449999999999999) -- (18.5,-6.449999999999999);
\draw[gray!45,line width=.35pt] (0,-7.6) -- (18.5,-7.6);
\draw[gray!45,line width=.35pt] (0,-8.75) -- (18.5,-8.75);
\draw[gray!45,line width=.35pt] (0,-9.899999999999999) -- (18.5,-9.899999999999999);
\draw[gray!45,line width=.35pt] (0,-11.049999999999999) -- (18.5,-11.049999999999999);
\node[anchor=north west,inner sep=0pt,align=left,text width=18.6cm,font=\fontsize{14}{17.78}\selectfont] at (0,-14) {\textbf{Problem 4:} A target is (a, b), with a and b positive whole numbers. Explain why a common divisor larger than 1 always produces a blocker. Does every blocker produce a common divisor?};
\draw[gray!45,line width=.35pt] (0,-17.0) -- (18.5,-17.0);
\draw[gray!45,line width=.35pt] (0,-18.15) -- (18.5,-18.15);
\draw[gray!45,line width=.35pt] (0,-19.3) -- (18.5,-19.3);
\draw[gray!45,line width=.35pt] (0,-20.45) -- (18.5,-20.45);
\draw[gray!45,line width=.35pt] (0,-21.6) -- (18.5,-21.6);
\draw[gray!45,line width=.35pt] (0,-22.75) -- (18.5,-22.75);\end{scope}\end{scope}\end{tikzpicture}
\newpage
\null\begin{tikzpicture}[remember picture,overlay,x=1cm,y=1cm]
\begin{scope}[shift={(current page.north west)}]
\node[anchor=north west,inner sep=0pt,font=\fontsize{11.5}{13}\selectfont] at (1.4,-.72) {Week 31 / Hidden orchard / Grades 4--5};
\node[anchor=south west,inner sep=0pt,font=\fontsize{9}{11}\selectfont] at (1.4,-27.2) {Bellingham Math Circle / Week 31 / F31-45-v2};
\node[anchor=south east,inner sep=0pt,font=\fontsize{9}{11}\selectfont] at (20.15,-27.2) {4};
\begin{scope}[shift={(1.4,-1.65)}]
\node[anchor=north west,inner sep=0pt,align=left,text width=18.6cm,font=\fontsize{14}{17.78}\selectfont] at (0,0) {\textbf{Problem 5:} List all the targets (a, b) whose two coordinates are each from 1 through 6 that have the same first visible dot as (6, 4). Repeat for (6, 3) and (6, 6). Explain why no target can belong to two different groups.};
\fill (3,-15.5) circle (1.3pt);
\fill (3,-13.5) circle (1.3pt);
\fill (3,-11.5) circle (1.3pt);
\fill (3,-9.5) circle (1.3pt);
\fill (3,-7.5) circle (1.3pt);
\fill (3,-5.5) circle (1.3pt);
\fill (3,-3.5) circle (1.3pt);
\fill (5,-15.5) circle (1.3pt);
\fill (5,-13.5) circle (1.3pt);
\fill (5,-11.5) circle (1.3pt);
\fill (5,-9.5) circle (1.3pt);
\fill (5,-7.5) circle (1.3pt);
\fill (5,-5.5) circle (1.3pt);
\fill (5,-3.5) circle (1.3pt);
\fill (7,-15.5) circle (1.3pt);
\fill (7,-13.5) circle (1.3pt);
\fill (7,-11.5) circle (1.3pt);
\fill (7,-9.5) circle (1.3pt);
\fill (7,-7.5) circle (1.3pt);
\fill (7,-5.5) circle (1.3pt);
\fill (7,-3.5) circle (1.3pt);
\fill (9,-15.5) circle (1.3pt);
\fill (9,-13.5) circle (1.3pt);
\fill (9,-11.5) circle (1.3pt);
\fill (9,-9.5) circle (1.3pt);
\fill (9,-7.5) circle (1.3pt);
\fill (9,-5.5) circle (1.3pt);
\fill (9,-3.5) circle (1.3pt);
\fill (11,-15.5) circle (1.3pt);
\fill (11,-13.5) circle (1.3pt);
\fill (11,-11.5) circle (1.3pt);
\fill (11,-9.5) circle (1.3pt);
\fill (11,-7.5) circle (1.3pt);
\fill (11,-5.5) circle (1.3pt);
\fill (11,-3.5) circle (1.3pt);
\fill (13,-15.5) circle (1.3pt);
\fill (13,-13.5) circle (1.3pt);
\fill (13,-11.5) circle (1.3pt);
\fill (13,-9.5) circle (1.3pt);
\fill (13,-7.5) circle (1.3pt);
\fill (13,-5.5) circle (1.3pt);
\fill (13,-3.5) circle (1.3pt);
\fill (15,-15.5) circle (1.3pt);
\fill (15,-13.5) circle (1.3pt);
\fill (15,-11.5) circle (1.3pt);
\fill (15,-9.5) circle (1.3pt);
\fill (15,-7.5) circle (1.3pt);
\fill (15,-5.5) circle (1.3pt);
\fill (15,-3.5) circle (1.3pt);
\node[font=\fontsize{10}{12}\selectfont,inner sep=1pt] at (3,-15.98) {0};
\node[font=\fontsize{10}{12}\selectfont,inner sep=1pt] at (5,-15.98) {1};
\node[font=\fontsize{10}{12}\selectfont,inner sep=1pt] at (7,-15.98) {2};
\node[font=\fontsize{10}{12}\selectfont,inner sep=1pt] at (9,-15.98) {3};
\node[font=\fontsize{10}{12}\selectfont,inner sep=1pt] at (11,-15.98) {4};
\node[font=\fontsize{10}{12}\selectfont,inner sep=1pt] at (13,-15.98) {5};
\node[font=\fontsize{10}{12}\selectfont,inner sep=1pt] at (15,-15.98) {6};
\node[font=\fontsize{10}{12}\selectfont,inner sep=1pt] at (2.5,-13.5) {1};
\node[font=\fontsize{10}{12}\selectfont,inner sep=1pt] at (2.5,-11.5) {2};
\node[font=\fontsize{10}{12}\selectfont,inner sep=1pt] at (2.5,-9.5) {3};
\node[font=\fontsize{10}{12}\selectfont,inner sep=1pt] at (2.5,-7.5) {4};
\node[font=\fontsize{10}{12}\selectfont,inner sep=1pt] at (2.5,-5.5) {5};
\node[font=\fontsize{10}{12}\selectfont,inner sep=1pt] at (2.5,-3.5) {6};
\node[font=\fontsize{12}{14}\selectfont,inner sep=1pt] at (2.55,-15.55) {O};
\draw[gray!45,line width=.35pt] (0,-19.0) -- (18.5,-19.0);
\draw[gray!45,line width=.35pt] (0,-20.15) -- (18.5,-20.15);
\draw[gray!45,line width=.35pt] (0,-21.3) -- (18.5,-21.3);
\draw[gray!45,line width=.35pt] (0,-22.45) -- (18.5,-22.45);\end{scope}\end{scope}\end{tikzpicture}
\newpage
\null\begin{tikzpicture}[remember picture,overlay,x=1cm,y=1cm]
\begin{scope}[shift={(current page.north west)}]
\node[anchor=north west,inner sep=0pt,font=\fontsize{11.5}{13}\selectfont] at (1.4,-.72) {Week 31 / Hidden orchard / Grades 4--5};
\node[anchor=south west,inner sep=0pt,font=\fontsize{9}{11}\selectfont] at (1.4,-27.2) {Bellingham Math Circle / Week 31 / F31-45-v2};
\node[anchor=south east,inner sep=0pt,font=\fontsize{9}{11}\selectfont] at (20.15,-27.2) {5};
\begin{scope}[shift={(1.4,-1.65)}]
\node[anchor=north west,inner sep=0pt,align=left,text width=18.6cm,font=\fontsize{14}{17.78}\selectfont] at (0,0) {\textbf{Problem 6:} Let d be the greatest whole number that divides both coordinates of a target. How does d determine the first dot after O and the number of blockers? Explain why there can be no extra dots between the ones your rule finds.};
\draw[gray!45,line width=.35pt] (0,-3.2) -- (18.5,-3.2);
\draw[gray!45,line width=.35pt] (0,-4.35) -- (18.5,-4.35);
\draw[gray!45,line width=.35pt] (0,-5.5) -- (18.5,-5.5);
\draw[gray!45,line width=.35pt] (0,-6.65) -- (18.5,-6.65);
\draw[gray!45,line width=.35pt] (0,-7.8) -- (18.5,-7.8);
\draw[gray!45,line width=.35pt] (0,-8.95) -- (18.5,-8.95);
\draw[gray!45,line width=.35pt] (0,-10.1) -- (18.5,-10.1);
\draw[gray!45,line width=.35pt] (0,-11.25) -- (18.5,-11.25);
\node[anchor=north west,inner sep=0pt,align=left,text width=18.6cm,font=\fontsize{14}{17.78}\selectfont] at (0,-14) {\textbf{Problem 7:} Find an entire row of visible dots, an infinite collection of hidden dots, and a target with exactly 99 blockers. Explain each choice.};
\draw[gray!45,line width=.35pt] (0,-17.0) -- (18.5,-17.0);
\draw[gray!45,line width=.35pt] (0,-18.15) -- (18.5,-18.15);
\draw[gray!45,line width=.35pt] (0,-19.3) -- (18.5,-19.3);
\draw[gray!45,line width=.35pt] (0,-20.45) -- (18.5,-20.45);
\draw[gray!45,line width=.35pt] (0,-21.6) -- (18.5,-21.6);
\draw[gray!45,line width=.35pt] (0,-22.75) -- (18.5,-22.75);\end{scope}\end{scope}\end{tikzpicture}
\end{document}
